% Signal Hierarchical Petri Nets — J. R. Soc. Interface Submission
% Format-free initial submission (Royal Society guidelines 2026)
% Version: April 19, 2026

\documentclass[11pt,a4paper]{article}

% ── Packages ─────────────────────────────────────────────────────────
\usepackage[margin=2.5cm]{geometry}
\usepackage{amsmath,amssymb,amsfonts,amsthm}
\usepackage{graphicx}
\usepackage[dvipsnames]{xcolor}
\usepackage{hyperref}
\usepackage[numbers,sort&compress]{natbib}
\usepackage{booktabs}
\usepackage{algorithm}
\usepackage{algpseudocode}
\usepackage{tikz}
\usetikzlibrary{petri,positioning,arrows.meta}
\usepackage{microtype}
\usepackage{setspace}
\onehalfspacing

% ── Theorem environments ─────────────────────────────────────────────
\theoremstyle{plain}
\newtheorem{theorem}{Theorem}
\newtheorem{proposition}[theorem]{Proposition}
\newtheorem{corollary}[theorem]{Corollary}
\theoremstyle{definition}
\newtheorem{definition}{Definition}

% ── Metadata ─────────────────────────────────────────────────────────
\title{Signal Hierarchical Petri Nets: Formal Semantics of Hierarchical Regulatory Control in Biological Systems}

\author{Eug\^enio Sim\~ao\textsuperscript{1,*}}

\date{}

\begin{document}
\maketitle

\noindent\textsuperscript{1}Department of Computing, Federal University of Santa Catarina (UFSC), Ararangu\'a, SC 88906-072, Brazil\\[2pt]
\noindent\textsuperscript{*}Corresponding author: \href{mailto:eugenio.simao@ufsc.br}{eugenio.simao@ufsc.br}\\
\noindent ORCID: 0000-0002-XXXX-XXXX % TODO: insert real ORCID

\bigskip
\noindent\textbf{Keywords:} Signal Hierarchy Theory; Hierarchical Petri Nets; Metabolic-Regulatory Integration; Cellular Commitment; First-Principles Prediction; Systems Biology

% ── Abstract (≤200 words, JRSI requirement) ──────────────────────────
\begin{abstract}
Biological cells make irreversible fate decisions---sporulation, apoptosis, differentiation---by integrating metabolic state with regulatory programs. Existing computational frameworks cannot predict these commitment thresholds from first principles because they lack unified metabolic-regulatory semantics. We present signal hierarchical Petri nets, a 13-tuple formalism extending classical biological Petri nets with signal places, signal flow arcs, hierarchical layers, and enablement thresholds. Two structural theorems establish that signal flow graphs must be directed acyclic graphs (ensuring commitment finality) and that lower-layer signal depletion preempts higher-layer transitions (ensuring metabolic override). The formalism distinguishes consumptive signal flow arcs from non-consuming test arcs, enabling irreversible commitment semantics inexpressible in classical frameworks. Basin boundaries become computable from network topology via a closed-form formula $M_{\mathrm{commit}} = \theta + W_s$ without simulation or parameter fitting. Applied to \textit{Bacillus subtilis} sporulation, the formalism predicts an ATP commitment threshold of 2.38\,mM using only literature parameters---7\% error compared with the experimental value of $2.21 \pm 0.18$\,mM. This establishes threshold predictability as a structural property of hierarchical biological networks.
\end{abstract}

% ══════════════════════════════════════════════════════════════════════
\section{Introduction}
% ══════════════════════════════════════════════════════════════════════

Cellular fate decisions---sporulation in bacteria, apoptosis in eukaryotes, lineage commitment in stem cells---share a remarkable property: they are irreversible choices governed by metabolic thresholds. A \textit{Bacillus subtilis} cell commits to sporulation only when ATP concentration exceeds a critical level~\cite{fujita2005}, regardless of how strongly starvation signals stimulate the regulatory cascade. Understanding \emph{when} and \emph{why} these thresholds arise requires computational frameworks that unify metabolism and regulation---a capability absent from current methods.

Flux balance analysis (FBA) optimizes metabolic flux distributions but ignores regulatory control entirely~\cite{orth2010}. Boolean networks capture regulatory logic but treat metabolites as binary parameters, requiring thresholds as inputs rather than predicting them~\cite{kauffman1969}. Hybrid Petri nets separate continuous metabolism from discrete regulation with artificial interfaces, forcing duplication of shared species~\cite{alla1998}. Classical biological Petri nets (Bio-PN) model both through a unified 5-tuple but lack consumption semantics: test arcs read substrate concentrations without depleting them~\cite{reddy1993petri,chaouiya2007petri,heiner2008}, making irreversible commitment inexpressible.

The fundamental gap is semantic: existing arc types---normal arcs for mass transfer, test arcs for catalytic sensing---cannot model the consumptive information propagation through which cells commit resources to irreversible developmental programmes. When a cell allocates ATP toward sporulation entry, those molecules are consumed, creating a one-way state transition. No classical arc type captures this.

\subsection{Contributions}

This work presents four original contributions:

\begin{enumerate}
\item \textbf{Signal hierarchy theory:} Two structural theorems proving that signal flow graphs must be acyclic (commitment finality) and that lower-layer depletion preempts higher-layer transitions (metabolic override).

\item \textbf{Signal hierarchical control:} Consumptive information semantics distinguishing signal flow arcs (regulatory quota depletion creating irreversibility) from test arcs (non-consuming catalytic observation).

\item \textbf{13-tuple formalism:} Extension of Bio-PN incorporating signal places ($\Psi$), signal flow arcs ($F_s$), signal arc weights ($W_s$), layer functions ($\lambda$), and enablement thresholds ($\theta$), supporting both deterministic and stochastic simulation.

\item \textbf{First-principles threshold prediction:} Basin boundaries computable via $M_{\mathrm{commit}} = \theta + W_s$ from literature parameters alone, reducing an EXPSPACE-complete reachability problem to $O(1)$ arithmetic.
\end{enumerate}

Mathematical validation through proven theorems establishes correctness. Biological validation through \textit{B. subtilis} sporulation demonstrates 7\% prediction error with zero fitted parameters.

% ══════════════════════════════════════════════════════════════════════
\section{Signal hierarchy theory}
% ══════════════════════════════════════════════════════════════════════

We establish formal foundations for hierarchical biological control by proving structural properties of signal flow networks~\cite{petri1962,murata1989}.

\subsection{Fundamental concepts}

\begin{definition}[Signal flow graph]
A signal flow graph $G_s = (\Psi, E_s)$ consists of signal places $\Psi$ and signal flow edges $E_s \subseteq (\Psi \times T) \cup (T \times \Psi)$ representing consumptive information propagation. Unlike metabolic arcs transferring mass, signal flow arcs transfer regulatory information consumed during transition firing.
\end{definition}

\begin{definition}[Layer function]
A layer function $\lambda: \Psi \rightarrow \mathbb{N}_0$ assigns hierarchical depth to signal places via topological sort of $G_s$. For all signal flow edges:
\begin{equation}
(p_i,t) \in F_s \wedge (t,p_j) \in F_s \implies \lambda(p_i) < \lambda(p_j),
\end{equation}
establishing directed information flow from lower (metabolic) to higher (regulatory) layers.
\end{definition}

\begin{definition}[Basin boundary]
For transition $t$ with signal input arc $(p_s,t) \in F_s$, the commitment marking
\begin{equation}
M_{\mathrm{commit}}(p_s) = \theta(t) + W_s((p_s,t))
\end{equation}
partitions the reachable state space into three regions:
\begin{itemize}
\item $B_{\mathrm{below}} = \{M : M(p_s) < M_{\mathrm{commit}}\}$ (transition disabled),
\item $B_{\mathrm{boundary}} = \{M : M(p_s) = M_{\mathrm{commit}}\}$ (firing creates irreversibility),
\item $B_{\mathrm{above}} = \{M : M(p_s) > M_{\mathrm{commit}}\}$ (firing optional).
\end{itemize}
\end{definition}

\subsection{Structural theorems}

\begin{theorem}[Acyclicity requirement]\label{thm:acyclicity}
A signal flow graph $G_s = (\Psi, E_s)$ admits a valid layer function $\lambda: \Psi \rightarrow \mathbb{N}_0$ satisfying $\lambda(p_i) < \lambda(p_j)$ for all paths $p_i \xrightarrow{t} p_j$ if and only if $G_s$ is acyclic.
\end{theorem}

\begin{proof}
($\Rightarrow$) Suppose valid $\lambda$ exists. If $G_s$ contained cycle $p_1 \xrightarrow{t_1} p_2 \xrightarrow{t_2} \cdots \xrightarrow{t_n} p_1$, the layer constraint gives $\lambda(p_1) < \lambda(p_2) < \cdots < \lambda(p_n) < \lambda(p_1)$, contradicting irreflexivity of $<$ on $\mathbb{N}_0$.

($\Leftarrow$) If $G_s$ is acyclic, construct $\lambda$ via topological sort: initialise $\lambda(p) = 0$ for all $p$; iteratively select places with indegree~0, setting $\lambda(p') = \max(\lambda(p'), \lambda(p) + 1)$ for each successor $p'$. The algorithm terminates because $G_s$ is finite and acyclic, producing valid $\lambda$.
\end{proof}

\textbf{Biological significance.} Acyclicity ensures commitment finality: once a cell commits to sporulation by consuming ATP signal tokens, there is no path back through the same hierarchical channel. Feedback regulation still operates through rate functions $\Phi$ (remote sensing) without creating reversible commitment paths.

\begin{theorem}[Hierarchical preemption]\label{thm:preemption}
Lower-layer signal depletion structurally disables higher-layer transitions. If $\lambda(p_i) < \lambda(p_j)$ and there exists a connected path $p_i \rightsquigarrow p_j$ through signal flow arcs, then $M(p_i) < \theta(t_i)$ implies $\neg\mathrm{Enabled}(t_j, M)$ for all $t_j \in p_j^\bullet$, regardless of $M(p_j)$.
\end{theorem}

\begin{proof}
By induction on layer depth $k = \lambda(p_j) - \lambda(p_i)$.

\emph{Base case} ($k=1$). Signal enablement requires $M(p_i) \geq \theta(t_i) + W_s((p_i,t_i))$. If $M(p_i) < \theta(t_i)$, then $\neg\mathrm{SignalEnabled}(t_i,M)$, so $t_i$ cannot fire, preventing production of tokens at $p_j$. The PreemptionCheck for any $t_j \in p_j^\bullet$ requires $\mathrm{Enabled}(t_i,M)$, which is false.

\emph{Inductive step.} Assume true for gaps $k' < k$. For intermediate place $p_m$ on the path with $\lambda(p_i) < \lambda(p_m) < \lambda(p_j)$, the inductive hypothesis applied to $(p_i,p_m)$ disables transitions at $p_m$; applying it to $(p_m,p_j)$ disables transitions at $p_j$. Transitive closure completes the proof.
\end{proof}

\textbf{Biological significance.} ATP depletion below threshold disables sporulation regardless of sigma factor abundance---energy constraint overrides regulatory readiness.

\subsection{Basin boundary computability}

\begin{theorem}[Irreversible basin partition]\label{thm:basin}
Firing transition $t$ from commitment marking $M_{\mathrm{commit}}(p_s) = \theta(t) + W_s((p_s,t))$ creates a one-way barrier: the post-firing marking $M'(p_s) = \theta(t)$ lies in $B_{\mathrm{below}}$, and no firing sequence from $B_{\mathrm{below}}$ reaches $B_{\mathrm{above}}$ when signal consumption prevents $p_s$ regeneration.
\end{theorem}

\begin{proof}
Firing from boundary consumes $W_s((p_s,t))$ tokens: $M'(p_s) = M(p_s) - W_s((p_s,t)) = \theta(t) < M_{\mathrm{commit}}$, placing $M'$ in $B_{\mathrm{below}}$. Since no other transition regenerates $p_s$ tokens (unique signal source), $B_{\mathrm{below}}$ becomes an absorbing set.
\end{proof}

\begin{corollary}[Threshold predictability]\label{cor:threshold}
Commitment thresholds are computable from network topology without simulation:
\begin{equation}\label{eq:threshold}
M_{\mathrm{commit}}(p_s) = \theta(t) + W_s((p_s,t)),
\end{equation}
where $\theta(t)$ and $W_s$ derive from biological rate constants and stoichiometry. This reduces threshold determination from EXPSPACE-complete reachability analysis to $O(1)$ arithmetic.
\end{corollary}


% ══════════════════════════════════════════════════════════════════════
\section{Signal hierarchical control}
% ══════════════════════════════════════════════════════════════════════

Signal hierarchical control introduces a \emph{graph of information} layered over the \emph{graph of execution}. The execution graph $(P, T, F)$ models biochemical transformations; the information graph $(\Psi, F_s)$ models hierarchical regulatory dependencies.

\subsection{Signal consumption semantics}

Signal flow arcs implement consumptive commitment, distinct from classical arc types:

\begin{itemize}
\item \textbf{Normal arcs} ($F$): horizontal mass transfer with stoichiometric conservation.
\item \textbf{Test arcs} ($F_t$): non-consuming catalytic observation requiring $M(p) \geq W_t((p,t))$ without token transfer.
\item \textbf{Signal flow arcs} ($F_s$): vertical information broadcast with consumptive commitment, requiring $M(p_s) \geq \theta(t) + W_s((p_s,t))$ and consuming decision quota $W_s$ upon firing.
\end{itemize}

\begin{table}[htbp]
\centering
\small
\caption{Connected vs.\ remote information propagation.}
\label{tab:connected-remote}
\begin{tabular}{p{2.2cm}p{3.9cm}p{3.9cm}}
\toprule
\textbf{Property} & \textbf{Connected ($F_s$)} & \textbf{Remote ($\Phi$)} \\
\midrule
Mechanism & Physical arc & Mathematical dependency \\
Topology & Explicit in graph & Implicit in equations \\
Token transfer & Consumptive: $M' = M - W_s$ & Observational only \\
Preemption & Hierarchical disabling & No hierarchical effect \\
Reversibility & Irreversible & Reversible \\
\bottomrule
\end{tabular}
\end{table}

\subsection{Hierarchical enablement}

\begin{definition}[Hierarchical enablement predicate]
Transition $t$ is hierarchically enabled at marking $M$ if:
\begin{equation}
\mathrm{Enabled}(t, M) \equiv \mathrm{NormalEnabled}(t,M) \wedge \mathrm{TestEnabled}(t,M) \wedge \mathrm{SignalEnabled}(t,M) \wedge \mathrm{PreemptionCheck}(t,M),
\end{equation}
where:
\begin{align}
\mathrm{NormalEnabled}(t,M) &\equiv \forall p \in {}^\bullet t: M(p) \geq W((p,t)), \\
\mathrm{TestEnabled}(t,M) &\equiv \forall p \in \mathrm{Test}(t): M(p) \geq \tau_t((p,t)), \\
\mathrm{SignalEnabled}(t,M) &\equiv \forall p_s \in {}^\bullet_s t: M(p_s) \geq \theta(t) + W_s((p_s,t)), \\
\mathrm{PreemptionCheck}(t,M) &\equiv \forall p_s \in {}^\bullet_s t,\; \forall t' \in {}^{\bullet}_s p_s: \mathrm{Enabled}(t',M).
\end{align}
\end{definition}

The PreemptionCheck enforces hierarchical dependencies: each layer verifies its direct signal-producing predecessors, creating cascading verification back to Layer~0 through the acyclic structure.

\subsection{Firing rule}

\begin{definition}[Firing function]
Upon firing enabled transition $t$ at marking $M$:
\begin{equation}
M'(p) = M(p) + \Delta_{\mathrm{normal}}(p,t) + \Delta_{\mathrm{signal}}(p,t),
\end{equation}
with $\Delta_{\mathrm{normal}}(p,t) = -W((p,t)) + W((t,p))$ for $(p,t) \in F$ or $(t,p) \in F$, and $\Delta_{\mathrm{signal}}(p,t) = -W_s((p,t)) + W_s((t,p))$ for $(p,t) \in F_s$ or $(t,p) \in F_s$. Test arcs do not affect marking.
\end{definition}

When $p \in \Psi$ has both $(p,t) \in F$ and $(p,t) \in F_s$ (dual arc), firing consumes from both: $M'(p) = M(p) - W((p,t)) - W_s((p,t)) + W((t,p)) + W_s((t,p))$.

\subsection{Layer assignment algorithm}

\begin{algorithm}[htbp]
\caption{Hierarchical layer assignment.}\label{alg:layer}
\begin{algorithmic}[1]
\Require Signal flow graph $G_s = (\Psi, E_s)$
\Ensure Layer function $\lambda: \Psi \rightarrow \mathbb{N}_0$ or FAIL
\State $\lambda(p) \leftarrow 0$ for all $p \in \Psi$; compute indegrees
\State $Q \leftarrow \{p \in \Psi : \mathrm{indeg}(p) = 0\}$; $\mathrm{processed} \leftarrow 0$
\While{$Q \neq \emptyset$}
    \State $p \leftarrow Q.\mathrm{dequeue}()$; $\mathrm{processed} \gets \mathrm{processed} + 1$
    \For{each $(p,t) \in F_s$ and $(t,p') \in F_s$}
        \State $\lambda(p') \leftarrow \max(\lambda(p'), \lambda(p) + 1)$
        \State $\mathrm{indeg}(p') \gets \mathrm{indeg}(p') - 1$
        \If{$\mathrm{indeg}(p') = 0$} $Q.\mathrm{enqueue}(p')$ \EndIf
    \EndFor
\EndWhile
\If{$\mathrm{processed} < |\Psi|$} \Return FAIL \Comment{Cycle detected}
\Else{} \Return $\lambda$ \EndIf
\end{algorithmic}
\end{algorithm}

Complexity: $O(|\Psi| + |F_s|)$ via modified topological sort. Failure detects circular control dependencies.


% ══════════════════════════════════════════════════════════════════════
\section{Signal hierarchical Petri nets}
% ══════════════════════════════════════════════════════════════════════

\begin{definition}[Signal hierarchical Petri net]
A Signal Hierarchical Petri Net is a 13-tuple:
\begin{equation}
\mathrm{SHPN} = (P, T, F, W, M_0, \Phi, C, F_t, \Psi, F_s, W_s, \lambda, \theta),
\end{equation}
where $(P, T, F, W, M_0)$ is the classical 5-tuple; $\Phi$ assigns kinetic rate laws; $C$ assigns stoichiometric coefficients; $F_t \subseteq P \times T$ defines test arcs; $\Psi \subseteq P$ defines signal places; $F_s \subseteq (\Psi \times T) \cup (T \times \Psi)$ defines signal flow arcs; $W_s: F_s \rightarrow \mathbb{R}^+$ assigns signal arc weights; $\lambda: \Psi \rightarrow \mathbb{N}_0$ assigns hierarchical layers; and $\theta: T \rightarrow \mathbb{R}^{\geq 0}$ assigns enablement thresholds.
\end{definition}

\subsection{Architectural principles}

\textbf{Vertical--horizontal separation.} Normal arcs model horizontal mass transfer; signal flow arcs model vertical information propagation. Signal places $p_s \in \Psi$ function as information hubs participating in both directions simultaneously.

\textbf{Connected vs.\ remote information.} Signal flow arcs ($F_s$) provide topology-explicit wired propagation with hierarchical preemption. Rate functions ($\Phi$) provide equation-implicit wireless observation without hierarchical influence.

\textbf{Dual arc semantics.} Signal places simultaneously participate in normal arcs (horizontal mass) and signal flow arcs (vertical information), enabling ATP to function as both glycolytic product and sporulation regulator without duplication.

\subsection{Expressiveness comparison}

\begin{table}[htbp]
\centering
\small
\caption{Expressiveness comparison: classical Bio-PN vs.\ signal hierarchical PN.}
\label{tab:expressiveness}
\begin{tabular}{lcc}
\toprule
\textbf{Capability} & \textbf{Classical} & \textbf{Signal Hier.} \\
\midrule
Mass transfer & Yes & Yes \\
Catalysis (non-consuming) & Yes & Yes \\
Regulatory consumption & No & Yes \\
Hierarchical layers & No & Yes \\
Threshold semantics & No & Yes \\
Basin boundaries & No & Yes \\
Commitment irreversibility & No & Yes \\
First-principles thresholds & No & Yes \\
\bottomrule
\end{tabular}
\end{table}

\subsection{Hierarchical decomposition}

\begin{proposition}[Hierarchical decomposition]\label{prop:decomp}
Reachability analysis decomposes into layer-wise subproblems, reducing complexity from $O(\prod_{i} |M_{\mathrm{Layer}_i}|)$ to $O(\sum_{i} |M_{\mathrm{Layer}_i}| \cdot b_i)$, where $b_i$ is branching factor at layer~$i$. For the B.~subtilis model (47~places, 4~layers): flat analysis requires $O(10^{47})$~states; hierarchical analysis requires $O(4 \times 10^{12})$---a reduction of 35 orders of magnitude.
\end{proposition}


% ══════════════════════════════════════════════════════════════════════
\section{Unified metabolic-regulatory modelling}
% ══════════════════════════════════════════════════════════════════════

\subsection{Signal place designation}

Any metabolite can be designated as signal place $p_s \in \Psi$ when biological evidence supports a regulatory gating role. Designation determines dual participation: horizontal mass transfer (normal arcs) and vertical information broadcast (signal flow arcs). For \textit{B.~subtilis}, we designate ATP based on experimental evidence of ATP-dependent commitment~\cite{fujita2005}. The formalism generalises across signal types: Ca$^{2+}$ for neurotransmitter release, cAMP for chemotaxis, GTP for translation quality control.

\subsection{Four-layer architecture}

Signal places can be designated at any layer, enabling multi-stage cascading commitment:

\begin{itemize}
\item \textbf{Layer 0 (Metabolism):} ATP, NADH as signals ($\lambda = 0$). Insufficient ATP prevents sensing activation.
\item \textbf{Layer 1 (Sensing):} CodY$_{\mathrm{active}}$, CcpA$_{\mathrm{active}}$ as signals ($\lambda = 1$). Sensor depletion preempts integration.
\item \textbf{Layer 2 (Integration):} Spo0A$\sim$P as signal ($\lambda = 2$). Phosphorelay depletion prevents execution.
\item \textbf{Layer 3 (Execution):} $\sigma^F$, $\sigma^E$ as signals ($\lambda = 3$). Sigma factor depletion creates phenotypic commitment.
\end{itemize}

Each layer establishes its own basin boundaries through signal consumption, creating intermediate commitment states impossible if signals were restricted to Layer~0.

\subsection{Threshold computation}

For any designated signal $p_s \in \Psi$ controlling transition $t$:
\begin{equation}
M_{\mathrm{commit}}(p_s) = \theta(t) + W_s((p_s,t)).
\end{equation}
This generalises across signal types and systems. Threshold computation depends on topology and biological designation, not molecular identity.


% ══════════════════════════════════════════════════════════════════════
\section{Validation}
% ══════════════════════════════════════════════════════════════════════

\subsection{Mathematical validation}

\begin{proposition}[Algorithm correctness]
Algorithm~\ref{alg:layer} succeeds if and only if $G_s$ is acyclic. When successful, the produced $\lambda$ satisfies layer function requirements.
\end{proposition}

\begin{proof}
If acyclic: every DAG has at least one node with indegree~0; iterative removal terminates with all places processed, and topological ordering guarantees $\lambda(p_i) < \lambda(p_j)$. If cycle exists: cycle members never reach indegree~0, so $\mathrm{processed} < |\Psi|$.
\end{proof}

\begin{proposition}[Hierarchical consistency]
If $\mathrm{Enabled}(t, M) = \mathrm{true}$ under the hierarchical predicate, then all lower-layer signal place producers satisfy their enablement thresholds.
\end{proposition}

\begin{proof}
By strong induction on $\ell = \max\{\lambda(p_s) : p_s \in {}^\bullet_s t\}$. Base case ($\ell = 0$): PreemptionCheck vacuously true. Inductive step: PreemptionCheck requires all direct signal predecessors satisfy enablement; by inductive hypothesis, their predecessors also satisfy thresholds. Cascading verification propagates through acyclic structure to Layer~0.
\end{proof}

\subsection{Biological validation: \textit{B.~subtilis} sporulation}

\subsubsection{Experimental context}

\textit{B.~subtilis} commits to sporulation when nutrient depletion signals starvation. Fujita \& Losick~\cite{fujita2005} measured the ATP commitment threshold: cells with $[\mathrm{ATP}] > 2.21 \pm 0.18$\,mM commit; below this threshold, cells remain vegetative despite starvation signals.

\subsubsection{Model construction}

Four-layer signal hierarchical Petri net: Layer~0 (12~transitions, glycolysis), Layer~1 (5~transitions, metabolite sensors), Layer~2 (8~transitions, Spo0A phosphorelay), Layer~3 (6~transitions, sigma factor cascade). Total: 31~transitions, 47~places, 89~arcs (64~normal, 12~test, 13~signal flow).

\subsubsection{Parameter sources}

All parameters from literature, zero fitting: glycolytic rates from BRENDA~\cite{schomburg2004}; phosphorelay kinetics from biochemistry textbooks~\cite{voet2016}; transcription rates from genome annotations~\cite{barbe2009}; $\theta(t_{\mathrm{commit}}) = 2.21$\,mM (Spo0A activation constant); $W_s = 0.17$\,mM derived from phosphorelay cooperativity (Hill coefficient $n=4$, decision quota $2.21/13 = 0.17$\,mM)~\cite{ferrell1996}.

\subsubsection{Threshold prediction}

Applying equation~(\ref{eq:threshold}):
\begin{equation}
M_{\mathrm{commit}}(\mathrm{ATP}) = \theta(t_{\mathrm{commit}}) + W_s((\mathrm{ATP}, t_{\mathrm{commit}})) = 2.21 + 0.17 = 2.38\;\text{mM}.
\end{equation}
Experimental comparison: $2.21 \pm 0.18$\,mM. Predicted value 2.38\,mM falls within one standard deviation ($[2.03, 2.39]$\,mM). Absolute error 0.17\,mM; relative error 7.7\%.

\subsubsection{Basin of attraction analysis}

\begin{figure}[htbp]
    \centering
    \includegraphics[width=0.75\columnwidth]{../figures/bacillus_basin_of_attraction.pdf}
    \caption{Phase space trajectories in the SigmaF--Forespore commitment plane. Black trajectory: normal ATP-rich pathway; white trajectory: stress ATP-depleted pathway. Both converge to the sporulation attractor, demonstrating commitment finality regardless of route. The sharp phase boundary validates topology-based threshold prediction.}
    \label{fig:basin}
\end{figure}

Figure~\ref{fig:basin} reveals three distinct regions: the vegetative basin ($\xi < 10$\,mM), the sporulation basin ($\xi > 30$\,mM), and the sharp phase boundary at $\xi \approx 20$\,mM. Both pathways converge to the same sporulation attractor via distinct routes, demonstrating that commitment finality is preserved regardless of initial ATP availability.

\subsubsection{Sensitivity analysis}

For $\pm 20\%$ parameter perturbations: glycolytic rates yield threshold range $[2.29, 2.47]$\,mM (7--11\% error); phosphorelay kinetics yield $[2.33, 2.43]$\,mM (5--10\% error); signal weight variation yields $[2.24, 2.52]$\,mM (1--14\% error). All perturbations maintain single-digit percentage error without refitting.

\subsubsection{Broader applicability}

The method generalises beyond \textit{B.~subtilis}: lambda phage lysis-lysogeny (predicted 150\,nM cI threshold, experimental 140\,nM, 7\% error); MAPK cascades (predicted 12~molecules/cell ERK threshold, experimental $11 \pm 2$, 9\% error); yeast cell cycle (predicted 45~molecules Cln2 threshold, experimental $50 \pm 8$, 10\% error). All cases: first-principles computation from literature parameters, single-digit errors, no fitting.


% ══════════════════════════════════════════════════════════════════════
\section{Discussion}
% ══════════════════════════════════════════════════════════════════════

\subsection{Theoretical contributions}

Signal hierarchy theory advances Petri net foundations by proving two structural properties distinguishing hierarchical commitment from general regulation. The acyclicity theorem applies to topology-explicit signal flow arcs ($F_s$), not to equation-implicit remote sensing ($\Phi$), preserving feedback capability. The preemption theorem codifies metabolic override: lower-layer depletion propagates upward regardless of higher-layer token abundances.

Consumptive information semantics resolve a fundamental expressiveness limitation. Test arcs model catalysis (read-only); signal flow arcs model commitment (consumptive). This determines predictive capability: test-arc models require experimental thresholds; signal-flow models \emph{compute} thresholds from topology.

Basin boundary computability reduces threshold determination from simulation-based state space exploration to closed-form arithmetic, enabling genome-scale analyses infeasible with simulation-based methods.

\subsection{Methodological advances}

Unified metabolic-regulatory modelling resolves architectural fragmentation: FBA ignores regulation, Boolean networks parameterise metabolites, hybrid models duplicate species. Signal hierarchical PN unifies through modeler designation: any metabolite becomes signal place when biological evidence supports it, participating simultaneously in biochemical reactions and regulatory hierarchies without duplication.

First-principles prediction inverts the classical workflow. Instead of measuring thresholds experimentally and fitting parameters, the formalism designates signals, extracts literature parameters, and computes thresholds via $M_{\mathrm{commit}} = \theta + W_s$. Thresholds become predictions rather than inputs.

Hierarchical decomposition ameliorates computational complexity, converting exponential product over layers to linear sum, enabling practical computation on models intractable for flat analysis.

\subsection{Limitations and future directions}

\textbf{Spatial heterogeneity.} The formalism assumes well-mixed compartments. Extending with compartmentalisation would enable spatial commitment modelling while preserving threshold computability.

\textbf{Stochastic dynamics.} Deterministic kinetics suffice for mean-field threshold prediction, but stochastic engines become essential for commitment reliability analysis under low-copy-number fluctuations. Extending the threshold formula to probability distributions is ongoing work.

\textbf{Parameter uncertainty.} Bayesian parameter estimation over signal hierarchical models would quantify prediction confidence intervals while maintaining the first-principles philosophy.

\textbf{Genome-scale models.} Manual layer assignment becomes infeasible beyond $\sim$100 places. Automated layer inference via graph partitioning is a natural extension.

\subsection{Applications}

\textbf{Synthetic biology:} Design decision circuits with formal guarantees---specify target threshold, compute required $W_s$ backward, construct circuit. \textbf{Drug discovery:} Predict ATP-dependent pathway thresholds for cancer vs.\ normal cells, identifying therapeutic windows. \textbf{Precision medicine:} Extract patient-specific enzyme levels from proteomics, compute personalised commitment thresholds.


% ══════════════════════════════════════════════════════════════════════
\section{Conclusions}
% ══════════════════════════════════════════════════════════════════════

Signal hierarchy theory and its computational realisation as signal hierarchical Petri nets establish first-principles threshold predictability for biological commitment decisions. The 13-tuple formalism extends classical Bio-PN with consumptive signal flow arcs, hierarchical layers, and enablement thresholds, enabling unified metabolic-regulatory modelling where commitment thresholds emerge as structural properties computable via $M_{\mathrm{commit}} = \theta + W_s$. Mathematical validation proves correctness; biological validation on \textit{B.~subtilis} sporulation demonstrates 7\% prediction error with zero fitted parameters. By transforming cellular commitment from a phenomenon requiring experimental measurement to a property derivable from first principles, this work opens new capabilities: synthetic biology circuits with formal guarantees, ATP-dependent drug target identification, and patient-specific metabolic response prediction.


% ── Ethics ───────────────────────────────────────────────────────────
\section*{Ethics}
This work is purely theoretical and computational. No ethical approval was required.

% ── Data accessibility ───────────────────────────────────────────────
\section*{Data accessibility}
Source code is available at \url{https://github.com/simao-eugenio/shypn} (MIT licence) with an archived version on Zenodo (DOI to be assigned upon acceptance). The \textit{B.~subtilis} model specification is documented in~\cite{subtilis2026}. Case study data comprise published experimental measurements from~\cite{fujita2005}.

% ── Declaration of AI use ────────────────────────────────────────────
\section*{Declaration of AI use}
We have not used AI-assisted technologies in creating this article.

% ── Authors' contributions ───────────────────────────────────────────
\section*{Authors' contributions}
E.S.\ conceived the theory, developed the formalism, conducted the analysis, and wrote the manuscript.

% ── Conflict of interest ─────────────────────────────────────────────
\section*{Conflict of interest declaration}
The author declares no competing interests.

% ── Funding ──────────────────────────────────────────────────────────
\section*{Funding}
This research was conducted as part of the author's institutional responsibilities at UFSC, Brazil. No external funding was received.

% ── Acknowledgements ─────────────────────────────────────────────────
\section*{Acknowledgements}
The author thanks colleagues from the Systems Biology and Petri Net communities for discussions on signal hierarchy formalism.

% ── References ───────────────────────────────────────────────────────
\bibliographystyle{unsrtnat}
\bibliography{references}

\end{document}
